\section{Introducción}
La navegación y planificación de rutas de costo mínimo sobre superficies topográficas reales representa un problema clásico en optimización. Ante la disponibilidad de mapas completos y entornos deterministas, algoritmos como el de Dijkstra garantizan la obtención del camino óptimo global. No obstante, un agente autónomo real opera bajo observación parcial y debe tomar decisiones locales con información incompleta. En tales escenarios, el Aprendizaje por Refuerzo (Reinforcement Learning, RL) provee un marco formal para el aprendizaje de políticas de navegación mediante interacción directa con el entorno.

En la fase previa del proyecto, el terreno fue modelado como un Proceso de Decisión de Markov (Markov Decision Process, MDP) y resuelto mediante Q-learning tabular en una cuadrícula reducida de $54\times54$ celdas con origen y destino fijos. Dicho enfoque almacena una matriz de valores y presenta severas limitaciones de escalabilidad y generalización, restringiendo la adaptabilidad del agente a condiciones o puntos de inicio no observados durante el entrenamiento.

En el presente estudio se incrementa la complejidad del escenario y se aplica Aprendizaje por Refuerzo Profundo (Deep Reinforcement Learning, DRL). Las principales contribuciones de este trabajo son:
\begin{itemize}
  \item Un entorno de mayor complejidad caracterizado por una resolución topográfica completa ($108\times108$ celdas), observación parcial (un parche local de $9\times9$ celdas y el vector relativo hacia la meta) e inicio aleatorio, lo cual demanda capacidades de generalización en el agente y sustenta la transición hacia aproximadores de función neuronales.
  \item Un análisis comparativo entre métodos basados en valor (Deep Q-Network y Double DQN) y métodos basados en política (REINFORCE) sobre el mismo terreno físico, empleando el algoritmo de Dijkstra como línea base de optimalidad.
  \item La caracterización experimental y resolución de dos problemáticas de diseño: la trampa de incentivos inducida por el rediseño de recompensa (*reward shaping*) potencial bajo descuento temporal y la inestabilidad/divergencia de los valores $Q$ en entornos con metas aleatorias.
\end{itemize}

En cuanto a la metodología, el terreno real se modela como un MDP a partir de un Modelo Digital de Elevación (Digital Elevation Model, DEM), donde el costo de tránsito es proporcional a la pendiente local. El agente recibe una observación en forma de parche local y el vector relativo al objetivo, aproximando la función $Q(s,a)$ mediante un perceptrón multicapa optimizado con DQN, memoria de reproducción de experiencias y red objetivo. El entrenamiento se realiza bajo un esquema de aprendizaje curricular (*curriculum learning*) incrementando progresivamente la distancia del estado inicial. La evaluación de la política codiciosa (*greedy*) se lleva a cabo sobre 20 configuraciones de inicio no observadas durante el entrenamiento, cuantificando la tasa de éxito y la brecha de costo en comparación con la ruta óptima de Dijkstra.

El resto del documento se estructura de la siguiente manera. La Sec.~\ref{sec:related} revisa la literatura científica previa y contextualiza la contribución de este trabajo. La Sec.~\ref{sec:teoria} describe el fundamento matemático de los algoritmos seleccionados. La Sec.~\ref{sec:metodo} formaliza el MDP, el espacio de observaciones y la estructura de la recompensa. La Sec.~\ref{sec:impl} detalla los aspectos de implementación, el pseudocódigo y los hiperparámetros. La Sec.~\ref{sec:exp} presenta los resultados cuantitativos y la Sec.~\ref{sec:disc} discute sus implicancias. Finalmente, la Sec.~\ref{sec:conc} presenta las conclusiones y líneas de trabajo futuro.
